\documentclass[lettersize,journal]{IEEEtran}
\usepackage{amsmath, amsfonts, amssymb, amstext}
\usepackage{physics}
\usepackage{siunitx}
\DeclareMathOperator*{\argmax}{arg\,max}
\DeclareMathOperator*{\argmin}{arg\,min}
\usepackage{graphicx}
\usepackage{xcolor}
\usepackage[font=footnotesize, labelsep=period, justification=raggedright, singlelinecheck=false]{caption}
\usepackage{subcaption}
\usepackage{array}
\usepackage{booktabs}
\usepackage{tabularx}
\usepackage{multirow}
\usepackage{algorithm}
\usepackage{algpseudocode}
\renewcommand{\algorithmicrequire}{\textbf{Input:}}
\renewcommand{\algorithmicensure}{\textbf{Output:}}
\usepackage{url}
\usepackage{verbatim}
\usepackage{cite}
\usepackage{stfloats}
\usepackage{balance}
\usepackage{soul}
\hyphenation{op-tical net-works semi-conduc-tor IEEE-Xplore}
\newcommand{\hlc}[2][yellow]{{\sethlcolor{#1} \hl{#2}}}
\colorlet{soulgreen}{green!30}
\newcommand\slee[1]{\hlc[soulgreen]{SL: -- #1 --}}
\colorlet{soulblue}{blue!30}
\newcommand\hpark[1]{\hlc[soulblue]{HBP: -- #1 --}}
\newcommand\mcc[2]{\multicolumn{1}{>{}p{#1}|}{#2}}
\begin{document}
\title{Link-Aware Routing in Multi-Tier LEO Mega-Constellations}
\author{Yoonsoo Choi, Anna Cho, Chang Kyung Kim, and SuKyoung Lee,~\IEEEmembership{Member,~IEEE}
\thanks{The authors are with the Department of Computer Science and Engineering, Yonsei University, Seoul, Republic of Korea (e-mail: yoonsoo@yonsei.ac.kr; annacho6@yonsei.ac.kr; ckim48@yonsei.ac.kr; sklee@yonsei.ac.kr).}
}
\markboth{IEEE Wireless Communications Letters,~Vol.~XX, No.~X, March~2026}%
{Yoonsoo Choi, \MakeLowercase{\textit{et al.}}: Link-Aware Routing in Multi-Tier LEO Mega-Constellations}
\maketitle

\begin{abstract}
In multi-tier low-Earth orbit (LEO) mega-constellations, the mobility of satellites across different orbital altitudes leads to dynamic changes in network topology and inter-satellite link (ISL) states, including ISL duration and capacity. These changes often result in unstable connectivity and disrupted end-to-end data transmission. To address this issue, we formulate a routing optimization problem to determine the optimal ISL path between two end users by maximizing the \textcolor{blue}{average ISL utility}, considering both ISL duration and capacity. To solve this problem, we propose a \textcolor{blue}{two-stage routing algorithm that first prunes unstable ISLs using a graph neural network (GNN) and then selects optimal end-to-end paths on the pruned graph using a utility-based heuristic}. Simulation results demonstrate that the proposed algorithm achieves higher throughput and a lower packet loss rate compared to benchmark methods.
\end{abstract}

\begin{IEEEkeywords}
Multi-tier LEO mega-constellations, routing, graph neural networks, inter-satellite link (ISL), ISL utility, throughput.
\end{IEEEkeywords}

\vspace{-0.1cm}
\section{Introduction}\label{sec:introduction}
Low Earth orbit (LEO) satellite-ground integrated networks have emerged as a promising solution to provide global Internet access. In the aerospace industry, several commercial enterprises, such as SpaceX, Amazon, and OneWeb, have constructed LEO mega-constellations (LMCs) to ensure ubiquitous coverage and support massive user access with high data rates~\cite{wang2022mega},~\cite{cao2025dense}\textcolor{blue}{,~\cite{he2024direct},~\cite{he2022multi}}.
Since a single LEO satellite cannot provide continuous service between two terrestrial users due to its limited coverage and mobility, inter-satellite links (ISLs) are introduced to enable seamless long-distance communication~\cite{mao2024intelligent}.
To enable reliable communication through ISLs under dynamic topology changes caused by satellite mobility, it is essential to determine optimal routing paths that maintain connectivity between satellites.
Hence, routing methods for ISLs have been proposed, primarily focusing on routing in single-tier LMCs, where satellites are deployed at the same orbital altitude.
However, the rapid growth of satellite deployments has resulted in multi-tier LMCs with heterogeneous orbital altitudes, increasing the number of deployed satellites and the corresponding ISLs per satellite. This leads to a larger search space for finding optimal ISL routing paths~\cite{luan2024fast}.


To overcome these limitations, several studies have proposed routing methods for multi-tier LMCs~\cite{hu2022software},~\cite{yan2024logic},~\cite{huang2026dynamic}.
Hu \textit{et al.}~\cite{hu2022software} proposed a multi-tier rectilinear Steiner tree construction algorithm by extending spanning-graph and edge-substitution approaches to three dimensions.
Yan \textit{et al.}~\cite{yan2024logic} proposed scalable routing methods by partitioning multi-tier LMCs into multiple groups based on their altitudes and orbital positions.
Huang \textit{et al.}~\cite{huang2026dynamic} proposed a bidirectional search strategy that integrates customized cost functions for multi-tier LMCs, reducing the search space through effective pruning while preserving path optimality.


However, none of these works consider ISL states between satellites in different tiers, including ISL duration and capacity.
In multi-tier LMCs, communication between satellites at different orbital altitudes experiences frequent ISL disruptions due to differences in orbital velocities. For instance, satellites at lower altitudes move faster than those at higher altitudes, resulting in rapidly changing connectivity and shorter ISL durations across different tiers.
Moreover, differences in ISL capacity may lead routing algorithms to select certain ISLs more frequently, resulting in traffic imbalances, congestion, and reduced network throughput~\cite{han2023time}.
Therefore, it is essential to consider ISL states, including duration and capacity, to improve network throughput for stable transmission in multi-tier LMCs.


In this letter, we propose a link-aware routing algorithm for multi-tier LMCs that handles dynamic changes in ISL states to enhance overall network throughput. We define an ISL utility based on ISL duration and capacity, and formulate a routing optimization problem to find optimal end-to-end (E2E) paths that maximize the \textcolor{blue}{average ISL utility}. As the formulated problem is NP-hard due to the exponential search space of all possible ISL combinations across source-destination pairs~\cite{guo2024lightweight}, we propose a routing algorithm to reduce the search space prior to path selection.
First, a graph neural network (GNN) prunes unstable ISLs based on predicted retention probability to construct a pruned graph. Subsequently, a heuristic algorithm determines the optimal E2E paths on the pruned graph. Simulation results demonstrate that the proposed algorithm outperforms existing approaches in terms of packet loss rate and throughput.

\vspace{-0.1cm}
\section{System Model and Problem Formulation}

\subsection{Multi-tier LMC System Model}

We consider a multi-tier LMC system consisting of a set of satellites $\mathcal{S}$, local controllers at ground stations, and a centralized LMC controller (LMCC) operating on a cloud server~\cite{hu2022software}. The LMC network is organized into multiple tiers $\mathcal{K}$, where each tier $k \in \mathcal{K}$ comprises orbits at the same altitude. Each satellite $i \in \mathcal{S}$ in tier $k$ moves along its respective orbital plane. Let $\mathcal{N}_i$ denote the set of satellites directly connected to satellite $i$ via intra-tier or inter-tier ISLs, where $j \in \mathcal{N}_i$ if the angular separation $\phi_{ij} \le \phi_{\mathrm{beam}}$ and $\phi_{\mathrm{beam}}$ is the beamwidth of the satellite antenna. The neighbor set $\mathcal{N}_i$ is further categorized into an intra-tier set $\mathcal{N}_i^{\mathrm{intra}}$ and an inter-tier set $\mathcal{N}_i^{\mathrm{inter}}$ based on their respective orbital altitudes. Let $\mathcal{R}$ denote the set of source--destination pairs served by the LMC network.
The LMCC maintains time-varying network snapshots generated from the two-line element (TLE) datasets, which provide satellite altitude $h_i$, inclination $\alpha_i$, and phase offset $\mu_i$. We assume that the network topology remains static within each snapshot interval $T$~\cite{han2023time}. \textcolor{blue}{However, ISL states including $l_{ij}$ and $c_{ij}$ are recomputed from the current satellite positions and velocities at each snapshot.}

As shown in Fig.~\ref{fig:system}, the procedure of the proposed algorithm at the beginning of each snapshot interval is as follows:
(1) Each satellite $i$ measures the signal-to-interference-plus-noise ratio (SINR) for its ISLs and transmits these measurements to the LMCC via its connected local controller.
(2) The LMCC calculates the duration ($l_{ij}$) and capacity ($c_{ij}$) of each ISL and computes the ISL utility ($U_{ij}$).
(3) The LMCC performs GNN-based link pruning on $\mathcal{G}$ to construct a pruned graph $\mathcal{G}'$.
(4) The LMCC determines routing paths ($P_{sd}^*$) on $\mathcal{G}'$ by a utility-based heuristic to maximize the \textcolor{blue}{average ISL utility}.
(5) Finally, the LMCC distributes the routing paths ($P_{sd}^*$) to the respective satellites through the local controllers.
\textcolor{blue}{The communication overhead is bounded by $|\mathcal{S}||\mathcal{N}_i|$ SINR reports per snapshot interval $T$, and each satellite receives only its own routing path $P_{sd}^*$ after path determination, so both uplink and downlink control overhead are negligible compared to $T$~\cite{han2023time},~\cite{hu2022software}.}

\begin{figure}[t]
\centering
\includegraphics[width=1\linewidth]{Image/SystemModel.PNG}
\caption{Overall procedure of the proposed algorithm.}
\label{fig:system}
\end{figure}

\subsubsection{ISL Duration}

Let $l_{ij}$ denote the ISL duration, defined as the time interval during which a link between satellites $i$ and $j$ remains available for data transmission~\cite{han2023time}. The position vector $\mathbf{p}_i$ of satellite $i$ is given by
\begin{equation} \label{eq:ecef_pos}
\mathbf{p}_i = (R + h_i) \left[\cos \varphi_i \cos \lambda_i,~ \cos \varphi_i \sin \lambda_i,~ \sin \varphi_i \right]^\top,
\end{equation}
where $R$ is the Earth's radius, and $\varphi_i$ and $\lambda_i$ are the geocentric latitude and longitude of satellite $i$, respectively. The ISL distance between satellites $i$ and $j$ is $d_{ij} = \|\mathbf{p}_j - \mathbf{p}_i\|$.

Let $\mathbf{v}_i$ and $\mathbf{v}_j$ denote the velocity vectors of satellites $i$ and $j$, respectively. The rate of change of the ISL distance $\dot{d}_{ij}$ is given by
\begin{equation} \label{eq:dist_rate}
\dot{d}_{ij} = \frac{\mathbf{v}_{ij} \cdot \mathbf{r}_{ij}}{\|\mathbf{r}_{ij}\|} = \|\mathbf{v}_{ij}\| \cos(\theta_{ij}),
\end{equation}
where $\mathbf{v}_{ij} = \mathbf{v}_j - \mathbf{v}_i$ is the relative velocity vector, $\mathbf{r}_{ij} = \mathbf{p}_j - \mathbf{p}_i$ is the relative position vector, and $\theta_{ij}$ is the angle between $\mathbf{v}_{ij}$ and $\mathbf{r}_{ij}$. Assuming a constant relative velocity during a short interval, the ISL duration $l_{ij}$ is given by
\begin{equation} \label{eq:link_duration}
l_{ij} = \begin{cases}
    \left| \frac{d_{ij}}{\dot{d}_{ij}} \right|, & \text{if } |\varphi_i|, |\varphi_j| \leq \varphi_p, \\
    0, & \text{otherwise},
\end{cases}
\end{equation}
where $\varphi_p$ is the polar latitude threshold.
\textcolor{blue}{Here, $|d_{ij}/\dot{d}_{ij}|$ represents the time until the link between satellites $i$ and $j$ becomes unavailable due to relative motion.}

\subsubsection{ISL Capacity}

The angular separation $\phi_{ij}$ is defined as the angle between the nadir vector of satellite $i$ and the ISL direction toward satellite $j$. Using the law of cosines, $\phi_{ij}$ is given by
\begin{equation} \label{eq:angular_separation}
\phi_{ij} = \arccos \left[\frac{(R+h_i)^2 + d_{ij}^2 - (R+h_j)^2}{2(R+h_i)d_{ij}}\right].
\end{equation}
When satellite $i$ transmits data to satellite $j$, the receiver $j$ experiences interference from other neighboring satellites~\cite{cao2025dense}. The intra-tier interference $I_j^{\mathrm{intra}}$ at satellite $j$ is given by
\begin{equation} \label{eq:interference_intra}
I_j^{\mathrm{intra}} = \sum_{j^{\prime} \in \mathcal{N}_{j}^{\mathrm{intra}} \setminus \{i\}} \frac{P_o G_t G_r b_{j^{\prime}j}^2}{PL(d_{j^{\prime}j})},
\end{equation}
where $P_o$ is the transmit power, $G_t$ and $G_r$ are the transmitter and receiver antenna gains, $PL(d_{j^{\prime}j})$ is the free-space path loss, and $b_{j^{\prime}j}$ is the beam misalignment factor. The inter-tier interference $I_j^{\mathrm{inter}}$ is given by
\begin{equation} \label{eq:interference_inter}
I_j^{\mathrm{inter}} = \sum_{j^{\prime} \in \mathcal{N}_j^{\mathrm{inter}}\setminus \{i\}} \eta(h_j,h_{j^{\prime}}) \frac{P_o G_t G_r b_{j^{\prime}j}^2}{PL(d_{j^{\prime}j})},
\end{equation}
where $\eta(h_j,h_{j^{\prime}}) = \exp(-\kappa |h_j-h_{j^{\prime}}|)$ is the altitude-dependent attenuation factor and $\kappa$ is the attenuation coefficient. \textcolor{blue}{The factor $\eta$ is applied only to inter-tier interferers, which arrive from off-boresight directions where the antenna gain is reduced~\cite{cao2025dense}.} The SINR for the ISL from satellite $i$ to $j$ is given by~\cite{cao2025dense}
\begin{equation} \label{eq:sinr}
\mathrm{SINR}_{ij} = \frac{P_o G_t G_r b_{ij}^2}{PL(d_{ij})\left(\textcolor{blue}{N_0 B} + I_{j}^{\mathrm{intra}} + I_{j}^{\mathrm{inter}}\right)},
\end{equation}
where $N_0$ is the noise power spectral density\textcolor{blue}{, $B$ is the available bandwidth, and $N_0 B$ denotes the total noise power}. The ISL capacity $c_{ij}$ is given by
\begin{equation}\label{eq:link_capacity}
    c_{ij} = B \log_2(1 + \mathrm{SINR}_{ij}).
\end{equation}

\subsection{Link Utility Model and Problem Formulation}

We model the LMC network as a graph $\mathcal{G} = (\mathcal{S}, \mathcal{E})$, where $\mathcal{E}$ denotes the set of ISLs. The quality of each ISL $(i,j) \in \mathcal{E}$ is quantified by a duration utility $U_{ij}^l$ normalized by the snapshot interval $T$ and a capacity utility $U_{ij}^c$ normalized via a softmax function\textcolor{blue}{, which suppresses the outlier sensitivity of min-max normalization and provides a stable relative ranking of link capabilities across all ISLs~\cite{fu2024softact}}, given by
\begin{equation} \label{eq:utility_components}
    U_{ij}^l = \frac{l_{ij}}{T}, \quad U_{ij}^c = \frac{\exp(c_{ij})}{\sum_{(m,n) \in \mathcal{E}} \exp(c_{mn})}.
\end{equation}
The ISL utility $U_{ij}$ is then defined as
\begin{equation} \label{eq:utility_def}
U_{ij} = \beta \cdot U_{ij}^l + (1-\beta) \cdot U_{ij}^c,
\end{equation}
where $\beta \in [0,1]$ is a weighting factor. \textcolor{blue}{Since $l_{ij}$ represents the time until the link becomes unavailable~\cite{han2023time} and $c_{ij}$ is determined by $\mathrm{SINR}_{ij}$ via Shannon's formula in Eq.~\eqref{eq:link_capacity}, ISLs with higher $U_{ij}$ are less likely to fail and provide higher achievable rate, both of which directly improve throughput~\cite{wang2023reliability}.}

Let $x_{ij}^{sd} \in \{0, 1\}$ denote a binary variable indicating whether ISL $(i,j) \in \mathcal{E}$ is included in the routing path for source--destination pair $(s,d) \in \mathcal{R}$. The optimization problem $\mathcal{P}$ is formulated as follows:
\begin{subequations} \label{eq:problem_formulation}
\begin{align}
\mathcal{P}: \max_{\{x_{ij}^{sd}\}} & \sum_{(s,d) \in \mathcal{R}} \frac{\sum_{(i,j)\in\mathcal{E}} x_{ij}^{sd} U_{ij}}{\sum_{(i,j)\in\mathcal{E}} x_{ij}^{sd}} \label{eq:obj} \\
\text{s.t.} \quad & \sum_{j} x_{sj}^{sd} - \sum_{j} x_{js}^{sd} = 1, \quad \forall (s,d)\in\mathcal{R} \label{eq:c1} \\
& \sum_{j} x_{dj}^{sd} - \sum_{j} x_{jd}^{sd} = -1, \quad \forall (s,d)\in\mathcal{R} \label{eq:c2} \\
& \sum_{j} x_{ij}^{sd} - \sum_{j} x_{ji}^{sd} = 0, \quad \forall i \neq s,d, \forall (s,d)\in\mathcal{R}. \label{eq:c3}
\end{align}
\end{subequations}
Constraints~\eqref{eq:c1} and~\eqref{eq:c2} enforce the flow requirements at the source and destination satellites, respectively. Constraint~\eqref{eq:c3} is the flow conservation condition for intermediate satellites.

\section{Link-Aware Routing Algorithm}
\label{sec:proposed_scheme}

Problem $\mathcal{P}$ is an Integer Non-Linear Programming (INLP) problem due to the binary variables $x_{ij}^{sd}$ and the non-convex objective function. This problem is known to be NP-hard because its search space grows exponentially with the number of ISL combinations across all pairs in $\mathcal{R}$~\cite{guo2024lightweight}. Consequently, finding an optimal solution directly is \textcolor{blue}{not tractable for real-time optimization} in LMCs~\cite{huang2026dynamic}. To address this, we propose a link-aware routing algorithm that reduces the search space before path selection. Specifically, our GNN model predicts the retention probability $\hat{p}_{ij}$ to prune unstable ISLs and construct a pruned graph $\mathcal{G}'$. Our heuristic algorithm then determines the final routing paths $P_{sd}^*$ on $\mathcal{G}'$ that maximize the \textcolor{blue}{average ISL utility} defined in Eq.~\eqref{eq:obj}.

\subsection{GNN-Based Link Pruning}
\label{subsec:learning}

We use a GNN because its neighborhood aggregation mechanism explicitly \textcolor{blue}{models} multi-hop topological dependencies\textcolor{blue}{, which are difficult to handle with} conventional heuristic algorithms~\cite{rao2025dynamic}. \textcolor{blue}{Unlike reinforcement learning (RL)-based approaches that require sequential agent-environment interactions for training~\cite{rao2025dynamic}, the GNN uses supervised learning with LP-derived labels generated offline, which produces deterministic per-link predictions with low inference latency.} By aggregating state information from surrounding satellites, the GNN predicts the retention probability $\hat{p}_{ij}$ for each ISL. Consequently, unstable ISLs are pruned based on $\hat{p}_{ij}$ to construct the pruned graph $\mathcal{G}'$.


\subsubsection{Feature Construction and Message Propagation}

Let $\mathcal{E}_i = \{(u,v) \in \mathcal{E} : u,v \in \mathcal{N}_i\}$ denote the set of ISLs among the neighbors of satellite $i$. The clustering coefficient $\mathrm{clust}(i)$ and betweenness centrality $\mathrm{bc}(i)$ are defined as
\begin{equation}
    \mathrm{clust}(i) = \frac{2 |\mathcal{E}_i|}{|\mathcal{N}_i|(|\mathcal{N}_i|-1)}, \quad \mathrm{bc}(i) = \sum_{p \neq q,\, p \neq i,\, q \neq i} \frac{n_{pq}(i)}{n_{pq}},
\end{equation}
where $n_{pq}$ is the total number of shortest paths between nodes $p$ and $q$, and $n_{pq}(i)$ is the number of those paths passing through satellite $i$.
Let $\mathbf{u}_i^{(0)}$ and $\mathbf{e}_{ij}$ denote the initial satellite and link feature vectors, respectively, defined as
\begin{equation}
    \mathbf{u}_i^{(0)} = [\, |\mathcal{N}_i|,\, \mathrm{clust}(i),\, \mathrm{bc}(i),\, h_i \,]^\top, \quad \mathbf{e}_{ij} = [\, \tilde{d}_{ij},\, \tilde{l}_{ij} \,]^\top,
\end{equation}
where $|\mathcal{N}_i|$ denotes the degree of satellite $i$, $h_i$ denotes the orbital altitude, $\tilde{d}_{ij} = d_{ij}/d_{\max}$ is the normalized distance, and $\tilde{l}_{ij} = l_{ij}/T$ is the normalized ISL duration. Here, $d_{\max}$ denotes the maximum possible ISL distance within the network.
Let $\mathbf{m}_{ij}^{(r)}$ denote the intermediate message from neighbor $j$ to satellite $i$ at iteration $r \in \{0, 1, \dots, L-1\}$. To capture topological dependencies during the message-passing phase, $\mathbf{m}_{ij}^{(r)}$ combines the features of the neighboring satellite and the connecting ISL as
\begin{equation} \label{eq:message}
    \mathbf{m}_{ij}^{(r)} = \mathrm{ReLU} \Big( \mathbf{W}_M^{(r)} [\, \mathbf{u}_j^{(r)} \,\|\, \mathbf{e}_{ij} \,] + \mathbf{b}_M^{(r)} \Big),
\end{equation}
where $\mathbf{W}_M^{(r)}$ and $\mathbf{b}_M^{(r)}$ are the weight matrix and bias vector for message generation, respectively, $\|$ represents vector concatenation, and $L$ is the number of message-passing iterations.
The updated satellite embedding $\mathbf{u}_i^{(r+1)}$ is obtained by aggregating the messages as
\begin{equation} \label{eq:satellite_update}
    \mathbf{u}_i^{(r+1)} = \sigma \Bigg( \mathbf{W}_S^{(r)} \mathbf{u}_i^{(r)} + \frac{1}{|\mathcal{N}_i|} \sum_{j \in \mathcal{N}_i} \mathbf{m}_{ij}^{(r)} \Bigg),
\end{equation}
where $\mathbf{W}_S^{(r)}$ is the learnable weight matrix for self-state transformation to preserve the satellite's intrinsic topological features during message aggregation, and $\sigma(\cdot)$ is a non-linear activation function. After $L$ iterations, using $\mathbf{u}_i^{(L)}$ obtained from Eq.~\eqref{eq:satellite_update}, the retention probability $\hat{p}_{ij}$ for each ISL is computed via a multi-layer perceptron (MLP) as
\begin{equation} \label{eq:readout}
    \hat{p}_{ij} = \mathrm{Sigmoid} \Big( \mathbf{w}_p^\top \mathrm{ReLU} \big( \mathbf{W}_p [\, \mathbf{u}_i^{(L)} \,\|\, \mathbf{u}_j^{(L)} \,\|\, \mathbf{e}_{ij} \,] + \mathbf{b}_p \big) \Big),
\end{equation}
where $\mathbf{W}_p$, $\mathbf{b}_p$, and $\mathbf{w}_p^\top$ are the learnable parameters of the MLP, and $\mathrm{Sigmoid}(\cdot)$ bounds the probability $\hat{p}_{ij} \in (0,1)$.

\subsubsection{Training and Pruning Optimization}

To generate ground-truth labels for GNN training, we solve a relaxed LP formulation of $\mathcal{P}$, where the binary variables $x_{ij}^{sd}$ are relaxed to continuous variables $\omega_{ij}^{sd} \in [0, 1]$. $\mathcal{P}$ is reformulated as the following LP:
\begin{equation} \label{eq:lp_relaxation}
\begin{aligned}
\mathcal{P}_{\text{LP}}: \max_{\{\omega_{ij}^{sd}\}} & \sum_{(s,d) \in \mathcal{R}} \sum_{(i,j) \in \mathcal{E}} \omega_{ij}^{sd} U_{ij} \\
\text{s.t.} \quad & \eqref{eq:c1} \text{--} \eqref{eq:c3}, \quad \forall (s,d) \in \mathcal{R}, \\
& 0 \leq \omega_{ij}^{sd} \leq 1, \forall (i,j) \in \mathcal{E}.
\end{aligned}
\end{equation}
\textcolor{blue}{Since the constraint matrix of \eqref{eq:c1}--\eqref{eq:c3} is totally unimodular, the LP relaxation yields integer-valued optimal solutions, which are obtained via the Gurobi optimizer and used as ground truth for GNN training~\cite{zhao2022link}.} Using $\omega_{ij}^{sd*}$ from Eq.~\eqref{eq:lp_relaxation}, the expected ISL utilization rate $\mathcal{I}_{ij}$ is computed as
\begin{equation} \label{eq:utilization}
    \mathcal{I}_{ij} = \frac{1}{|\mathcal{R}|} \sum_{(s,d)\in\mathcal{R}} \omega_{ij}^{sd*},
\end{equation}
and the corresponding binary ground-truth label $y_{ij}$ is assigned using Eq.~\eqref{eq:utilization} as
\begin{equation} \label{eq:label}
    y_{ij} = \begin{cases} 1, & \text{if } \mathcal{I}_{ij} \geq \tau_{\text{train}} \\ 0, & \text{otherwise,} \end{cases}
\end{equation}
where $\tau_{\text{train}}$ denotes the labeling threshold for ISL retention.
Let $\boldsymbol{\Theta}$ denote the set of all trainable parameters in the GNN. We optimize $\boldsymbol{\Theta}$ via the binary cross-entropy loss $\mathcal{L}(\boldsymbol{\Theta})$:
\begin{equation} \label{eq:loss_function}
\begin{aligned}
    \mathcal{L}(\boldsymbol{\Theta}) = &-\frac{1}{|\mathcal{E}|} \sum_{(i,j) \in \mathcal{E}} \Big[ y_{ij} \log \hat{p}_{ij} \\
    &+ (1 - y_{ij}) \log (1 - \hat{p}_{ij}) \Big] + \Lambda_{wd} \|\boldsymbol{\Theta}\|^2,
\end{aligned}
\end{equation}
where $\Lambda_{wd}$ is the weight decay coefficient. After the training phase, the pruned graph $\mathcal{G}' = (\mathcal{S}, \mathcal{E}')$ is constructed with a reduced link set $\mathcal{E}' = \{ (i,j) \in \mathcal{E} \mid \hat{p}_{ij} \geq \tau \}$, where $\tau$ represents the inference pruning threshold. \textcolor{blue}{Since $\mathbf{e}_{ij}$ is recomputed from the current SINR and TLE data at each snapshot, the trained GNN is reused across snapshots within the same orbital configuration. Since $\mathbf{u}_i^{(0)}$ and $\mathbf{e}_{ij}$ are defined solely from structural and link state properties without dependence on the source--destination pairs in $\mathcal{R}$, the trained GNN is also applicable to different source--destination configurations within the same orbital configuration without retraining. The inference complexity is $\mathcal{O}(L|\mathcal{E}|d)$, where $d$ is the hidden dimension.}

\subsection{Utility-Based Heuristic Routing on the Pruned Graph}
\label{subsec:routing}

Since the pruned graph $\mathcal{G}'$ retains only stable and high-utility ISLs, we determine the routing paths $P_{sd}^*$ on $\mathcal{G}'$ using a priority-based heuristic. To evaluate both the accumulated path utility and the maximum possible utility from the next hop, the priority function $f(i)$ for satellite $i$ is defined as
\begin{equation} \label{eq:priority_func}
    f(i) = \frac{g(i) + \psi(i)}{H_{si} + 1},
\end{equation}
where $g(i)$ denotes the accumulated ISL utility from source satellite $s$ to current satellite $i$, $H_{si}$ denotes the corresponding hop count, and $\psi(i) = \max_{j \in \mathcal{N}_i \cap \mathcal{E}'} U_{ij}$ is the maximum ISL utility among all available outgoing ISLs connected to satellite $i$ within the pruned graph $\mathcal{G}'$. \textcolor{blue}{It is noted that $f(i)$ in Eq.~\eqref{eq:priority_func} balances path utility and hop count: it assigns a higher priority to paths whose average per-hop utility is high relative to their length, which penalizes paths that achieve shorter hop counts at the expense of low-utility ISLs.}

A priority queue $Q$ is initialized with the source satellite $s$. At each iteration, the algorithm extracts satellite $i$ with the highest priority $f(i)$ and evaluates its adjacent neighbors $j \in \mathcal{N}_i \cap \mathcal{E}'$. The candidate priority is computed as $f_n = (g(i) + U_{ij} + \psi(j)) / (H_{si} + 2)$. If $j$ is unvisited or $f_n > f(j)$, the routing state of $j$ is updated and $j$ is re-inserted into $Q$ with satellite $i$ recorded as its predecessor. This expansion continues until the destination $d$ is reached, after which the optimal path $P_{sd}^*$ is extracted by backtracking the recorded predecessors from $d$ to $s$. \textcolor{blue}{The priority queue expansion has complexity $\mathcal{O}(|\mathcal{S}|\log|\mathcal{S}| + |\mathcal{E}'|)$ per source--destination pair.}

\section{Performance Evaluation}

\subsection{Simulation Settings}

The simulation of the multi-tier LMC is conducted using the Starlink TLE dataset and the Satellite Tool Kit (STK) to \textcolor{blue}{model} realistic time-varying orbital dynamics~\cite{luan2024fast}. We consider 1,500 satellites uniformly distributed across $K=3$ tiers, where each tier comprises 500 satellites distributed across 17 orbital planes. We evaluate two altitude scenarios for the three tiers: (1) Scenario~1 with altitudes of 540, 550, and 560~km, and (2) Scenario~2 with altitudes of 400, 500, and 600~km.

The proposed algorithm is implemented using PyTorch 1.9.0 on a workstation equipped with an NVIDIA RTX 3080 Ti GPU. The GNN architecture consists of $L=3$ graph convolutional layers with a hidden dimension of $64$. The model is trained over $100$ epochs using the Adam optimizer with a learning rate of $10^{-3}$. The channel bandwidth $B$ is $500$~MHz, the snapshot interval $T$ is $5$~min, and the noise power spectral density $N_0$ is $-174$~dBm/Hz~\cite{han2023time}. \textcolor{blue}{Doppler-induced frequency offsets in LEO ISLs are compensated using standard frequency-tracking methods~\cite{darsena2024hybrid}, and are assumed to be well-managed in standard LEO ISL operation~\cite{cao2025dense}.} The parameters $\varphi_p$ and $P_o$ are set to $70^\circ$ and $10$~W, respectively. The weighting factor $\beta$ and the pruning threshold $\tau$ are both set to $0.5$\textcolor{blue}{, selected based on the sensitivity analysis detailed in Section~IV-B}. \textcolor{blue}{To verify optimality, we generate reduced 3-tier networks of $N \in \{30, 50, 70\}$ satellites with the same constellation configuration and enumerate all simple paths up to 8 hops to obtain the optimal solution.}

We evaluate the performance of the proposed algorithm in terms of the packet loss rate (PLR) and throughput compared with the following benchmark methods:
1) Distance-based Routing (DR)~\cite{wang2023reliability},~\cite{hu2022software}: the routing path is selected based on the geometric distance between satellites to minimize transmission delay;
2) Static Topology-based Routing (STR)~\cite{yan2024logic},~\cite{huang2026dynamic}: the routing decision is made by utilizing a fixed network structure within each snapshot interval;
\textcolor{blue}{3) Graph Reinforcement Learning Routing (GRLR)~\cite{zhang2025grlr}: the routing path is selected by combining GNNs and RL to minimize the end-to-end delay.}
\textcolor{blue}{Throughput is measured as the total successfully delivered data volume per unit time across all source--destination pairs, and PLR is the fraction of packets that fail to reach the destination due to link failures.}

\subsection{Simulation Results}

\begin{figure}[t!]
\centering
\includegraphics[width=0.85\linewidth]{Image/fig2_combined_duration.eps}
\vspace{-0.2cm}
\caption{Average throughput and packet loss rate over varying data transmission durations.}
\vspace{-0.2cm}
\label{fig:combined_results_duration}
\end{figure}

Fig.~\ref{fig:combined_results_duration} illustrates the average throughput and packet loss rate with respect to the data transmission duration at a data transmission rate of $2.8$~Gbps. In Scenario~1, the proposed algorithm achieves 3.0\% and 5.9\% higher throughput than DR and STR at a duration of 10~s, \textcolor{blue}{increasing to 5.4\%, 8.5\%, and 12.7\% higher throughput than GRLR, DR, and STR at 100~s, respectively}. This is because the proposed algorithm selects ISLs with longer durations and higher SINR values, as shown in Fig.~\ref{fig:box_comparison}, while \textcolor{blue}{GRLR does not consider ISL duration or capacity,} DR experiences increasing link failures over time, and STR experiences more disconnections as the fixed snapshot becomes outdated. Consequently, the packet loss rate is reduced by 2.1\% and 3.3\% at 10~s compared to DR and STR, \textcolor{blue}{and by 3.6\%, 5.7\%, and 7.4\% at 100~s compared to GRLR, DR, and STR}, respectively, as shown in Fig.~\ref{fig:combined_results_duration}(b). In Scenario~2, the proposed algorithm achieves 2.2\% and 5.6\% higher throughput than DR and STR at 10~s, \textcolor{blue}{increasing to 5.5\%, 8.1\%, and 12.8\% higher throughput than GRLR, DR, and STR at 100~s, respectively}, and reduces the packet loss rate by 1.2\% and 2.6\% at 10~s compared to DR and STR, \textcolor{blue}{and by 3.4\%, 5.2\%, and 7.1\% at 100~s compared to GRLR, DR, and STR}, respectively, as shown in Figs.~\ref{fig:combined_results_duration}(c) and (d).

\begin{figure}[t!]
\centering
\includegraphics[width=0.85\linewidth]{Image/fig3_combined_box.eps}
\vspace{-0.2cm}
\caption{Comparison of average ISL duration and SINR.}
\vspace{-0.2cm}
\label{fig:box_comparison}
\end{figure}

\begin{figure}[t!]
\centering
\includegraphics[width=0.85\linewidth]{Image/fig4_combined_rate.eps}
\vspace{-0.2cm}
\caption{Average throughput and packet loss rate over varying data transmission rates.}
\vspace{-0.2cm}
\label{fig:combined_results_rate}
\end{figure}

Fig.~\ref{fig:combined_results_rate} shows the average throughput and packet loss rate as the data transmission rate increases at a data transmission duration of $100$~s. In Scenario~1, the proposed algorithm achieves 6.7\% and 9.1\% higher throughput than DR and STR at 1.0~Gbps, \textcolor{blue}{increasing to 5.4\%, 8.5\%, and 12.7\% higher throughput than GRLR, DR, and STR at 2.8~Gbps, respectively}. This is because the proposed algorithm jointly considers ISL duration and capacity to select paths with stable ISLs, while \textcolor{blue}{GRLR does not consider ISL duration or capacity,} DR selects paths based on geometric distance without considering ISL capacity, and STR uses a fixed topology that does not reflect the current ISL state. Consequently, the packet loss rate is reduced by 5.6\% and 7.3\% at 1.0~Gbps compared to DR and STR, \textcolor{blue}{and by 3.6\%, 5.7\%, and 7.4\% at 2.8~Gbps compared to GRLR, DR, and STR}, respectively, as shown in Figs.~\ref{fig:combined_results_rate}(a) and (b). In Scenario~2, the proposed algorithm achieves 6.1\% and 8.7\% higher throughput than DR and STR at 1.0~Gbps, \textcolor{blue}{increasing to 5.5\%, 8.1\%, and 12.8\% higher throughput than GRLR, DR, and STR at 2.8~Gbps, respectively}, and reduces the packet loss rate by 5.1\% and 6.8\% at 1.0~Gbps compared to DR and STR, \textcolor{blue}{and by 3.4\%, 5.2\%, and 7.1\% at 2.8~Gbps compared to GRLR, DR, and STR}, respectively, as shown in Figs.~\ref{fig:combined_results_rate}(c) and (d).

\textcolor{blue}{Table~\ref{tab:sensitivity} shows the sensitivity of throughput and PLR to $\tau$, $\beta$, and $T$. We observe that $\tau = \beta = 0.5$ and $T = 5$~min achieve the best performance. This is because $\tau \geq 0.7$ removes too many ISLs and $T \geq 10$~min cannot reflect the current ISL state. The proposed algorithm also achieves 8.8\% higher throughput and 5.2\% lower PLR than its version without GNN-based pruning. Table~\ref{tab:optimality} confirms near-optimal performance, with the proposed algorithm reaching 95.7\% and 93.0\% of the optimal throughput at $N=30$ and $N=70$, respectively.}

\begin{table}[t]
\caption{Sensitivity of $\tau$, $\beta$, and Snapshot Interval $T$ (Scenario~2, 2.8~Gbps, 100~s; Throughput in Gbps, PLR in \%)}
\label{tab:sensitivity}
\centering
\footnotesize
\setlength{\tabcolsep}{3.5pt}
\begin{tabular}{cc|cc|cc}
\hline
$\tau$ & Thr./PLR & $\beta$ & Thr./PLR & $T$ (min) & Thr./PLR \\
\hline
0.1 & 1.461/26.0  & 0.1 & 1.529/23.6 & 1  & 1.589/20.8 \\
0.3 & 1.461/26.0  & 0.3 & 1.577/21.6 & 5  & 1.589/20.8 \\
0.5 & 1.589/20.8  & 0.5 & 1.589/20.8 & 10 & 0.863/57.0 \\
0.7 & 0.024/98.8  & 0.7 & 1.548/22.7 & 30 & 0.319/84.1 \\
0.9 & 0.000/100.0 & 0.9 & 1.510/24.4 & -- & --         \\
\hline
\end{tabular}
\end{table}

\begin{table}[t]
\caption{Optimality at Small Scale (Scenario~2, 2.8~Gbps, 100~s)}
\label{tab:optimality}
\centering
\footnotesize
\begin{tabular}{l|ccc}
\hline
$N$                              & 30          & 50          & 70          \\
\hline
Throughput Opt/Prop (Gbps)       & 2.151/2.060 & 2.058/1.906 & 2.123/1.974 \\
PLR Opt/Prop (\%)                & 6.3/7.2     & 7.9/9.9     & 6.7/8.3     \\
\hline
\end{tabular}
\end{table}

\section{Conclusion}
In this letter, we proposed a link-aware routing algorithm for multi-tier LMCs that maximizes the \textcolor{blue}{average ISL utility} via GNN-based ISL pruning and utility-based heuristic routing on the pruned graph. Simulation results demonstrate that the proposed algorithm outperforms the benchmark methods in terms of packet loss rate and throughput.

\bibliographystyle{IEEEtran}
\bibliography{ref}
\end{document}
